\section{Ableitungen in $\R^n$}

\Title{Polynome in $\R^n$}

Ein \textbf{Monom} von Grad $e$ ist eine Funktion
\begin{align*}
    (x_1, \dots, x_n) \mapsto {x_1}^{d_1} \cdot \ldots \cdot {x_n}^{d_n}\\
    e = d_1 + \ldots + d_n
\end{align*}

\begin{tcolorbox}
    Ein \textbf{Polynom} mit $n$ Variablen von Grad $\leq d$ ist eine endliche Summe von Monomen mit 
    Grad $e \leq d$.
\end{tcolorbox}


\Title{Konvergenz}

Ähnlich zu Funktionen in $\R$ können wir auch die Konvergenz für Funktionen $f$ in $\R^n$ definieren.
Aber zuerst repetieren wir ein paar Begriffe der Linearen Algebra:

\begin{enumerate}[label=(\arabic*)]
    \item Skalarprodukt: $<x,y> = \Sum_{i = 0} x_i \cdot y_i$
    \item Euklidische Norm: $||x|| := \sqrt{x_1^2 + \dots + x_n^2}$
\end{enumerate}

Die euklidische Norm hat folgende Eigenschaften:

\begin{enumerate}[label=(\arabic*)]
    \item $||x|| \geq 0$, $||x|| = 0$ genau dann wenn $x = 0$
    \item $||\lambda x|| = |\lambda| \cdot ||x||, \forall \lambda \in \R$
    \item $||x + y|| \leq ||x|| + ||y||$ ($\triangle$-Ungleichung)
    \item $|<x,y>| \leq ||x|| \cdot ||y||$
\end{enumerate}


%-------------------------
\vfill\null
\columnbreak
%-------------------------


Nun können wir \textbf{Konvergenz} definieren:
\begin{tcolorbox}
    Sei $(x_k)_{k \in \mathbb{N}}, x_k \in \R^n$. Die folgenden Definitionen sind äquivalent für $\lim_{k \to \infty} x_k = y$.

    \begin{enumerate}[label=(\arabic*)]
        \item $\forall \varepsilon > 0 \; \exists N \geq 1$ so dass $\forall k \geq N \quad ||x_k - y|| < \varepsilon$
        \item Für jedes $i$, $1 \leq i \leq n$ konvergiert die Folge $(x_{k, i})_k$ von reellen Zahlen nach $y_i$.
        \item Die Folge der reellen Zahlen $||x_k - y||$ konvergiert nach $0$.
    \end{enumerate}
\end{tcolorbox}

Sei $f: \X \subset \R^n \to \R^m$ und $x_0 \in \X, \; y \in \R^m$. Wir sagen $f$ hat den Limes $y$ 
für $x \to x_0$ wenn $x \neq x_0$ und etwas der Folgenden zutrifft:

\begin{enumerate}[label=(\arabic*)]
  \item Für alle $ \varepsilon > 0 \; \exists \delta > 0$ s.t. $\forall x \in \X,\; x \neq x_0$ so dass
  $||x - x_0|| < \delta$ gilt $||f(x) - y|| < \varepsilon$.
  \item $\forall$ Folgen $(x_k)$ in $\X$ mit $\lim x_k = x_0$ und $x_k \neq x_0$ konvergiert die Folge $f(x_k)$ nach $y$.
\end{enumerate}


\Title{Stetigkeit}

Sei $f: \X \subset \R^n \to \R^m$ und $x_0 \in \X$. Wir sagen $f$ ist stetig in $x_0$ falls einer der Folgenden
Bedingungen erfüllt ist:

\begin{enumerate}[label=(\arabic*)]
  \item $\forall \varepsilon > 0 \; \exists \delta > 0$ so dass $\forall x \in \X, \; ||x - x_0|| < \delta \Rightarrow$ 
  $||f(x) - f(x_0)|| < \varepsilon$.
  \item $\forall$ Folgen $(x_k)$ in $\X$ mit $\lim x_k = x_0$ gilt $\lim f(x_k) = f (\lim x_k)$.
\end{enumerate}

$f$ ist stetig in $\X$ falls $f$ für jeden Punkt $x_0 \in \X$ stetig ist. Folgende Eigenschaften gelten.

\begin{enumerate}[label=(\arabic*)]
  \item $f(x = x_1, \dots, x_n) \mapsto (f_1(x), \dots, f_m(x))$ und $f_i: \R^n \mapsto \R$,  
  dann gilt $f: \R^n \mapsto \R^m$ stetig $\Leftrightarrow f_i \; \forall i = 1, \dots, m$ sind 
  stetig.
  \item Lineare Funktionen $x \mapsto Ax$ und Polynome sind stetig.
  \item Summen und Produkte von stetigen Funktionen sind stetig.
  \item Funktionen von unterschiedlichen Variablen sind stetig, falls alle Variablen stetig sind.
  \item Verknüpfungen von stetigen Funktionen sind stetig.
\end{enumerate}


\Title{Sandwich Lemma}

Wenn $f, g, h: \R^n \to \R$ Funktionen sind mit $f(x) < g(x) < h(x), \forall x \in \R^n$, so gilt:
$$\lim_{x \to a} f(x) = \lim_{x \to a} h(x) = L \Rightarrow \lim_{x \to a} g(x) = L$$


\Title{Definition von Mengen}

Eine Menge $\X \subset \R^n$ ist \textbf{beschränkt}, falls die Menge $\{ ||x|| \;|\; x \in \X\}$ beschränkt 
ist in $\R$.
$$\exists R \geq 0, \forall x \in \X : ||x|| \leq R$$

Eine Menge $\X \subset \R^n$ ist \textbf{abgeschlossen}, falls jede Folge $(x_k)_{k \in \mathbb{N}} 
\subset \X$ die in $\R^n$ konvergiert, zu einem Punkt in $y \in \X$ konvergiert. Hier ist es hilfreich sich
eine Ball vorzustellen. Für Gegenbeispiele wird oft $\frac{1}{k}$ oder $<$ gebraucht. \medskip

Eine Menge ist \textbf{kompakt}, falls sie beschränkt und abgeschlossen ist.

\begin{tcolorbox}
    $f: \R^n \to \R^m$ ist genau dann stetig, falls das Urbild jeder abgeschlossenen Menge $\subset \R^m$ stetig ist.
    Achtung: Die impliziert weder Kompaktheit noch Beschränktheit!
\end{tcolorbox}

Eine Menge $\X \subset \R^n$ ist \textbf{offen}, wenn ihr Komplement $\R^n \setminus \X$ abgeschlossen ist.
Dies ist äquivalent zu $\forall x \in \X, \; \exists r > 0$ so dass
$\{y \in \R^n \;|\; ||y - x|| < r\} = B_r(x) \subset \X$. \medskip


%-------------------------
\vfill\null
\columnbreak
%-------------------------


Hier einige Beispiele:

\begin{enumerate}[label=(\arabic*)]
  \item $(a, b) \subset \R$ ist offen.
  \item $[a, b) \subset \R$ ist weder offen noch abgeschlossen.
  \item $\R^n$ und $\emptyset$ sind beide offen und abgeschlossen.
  \item $(a_1, b_1) \times (a_2, b_2) \subset \R^2$ ist offen.
  \item Urbilder von von offenen Mengen sind unter stetigen Abbildungen offen (gilt auch für abgeschlossen).
\end{enumerate}

Eine Menge $\X \subset \R^n$ ist \textbf{Konvex}, falls $\forall x,y \in \X: \lambda x + (1-\lambda)y \in \X$ gilt.


\Title{Bolzano-Weierstrass}

Jede beschränkte Folge in $\R^n$ hat eine konvergente Teilfolge.


\Title{Min-Max Theorem}

Sei $\X \subset \R^n, \X \neq \emptyset$ eine kompakte Menge und $f: \X \to \R$ eine stetige
Funktion. Dann ist $f$ beschränkt und ein globales Maximum $x^+$ und ein globales Minimum $x^-$
existieren, so dass

$$f(x^+) = \sup_{x \in \X} f(x), \quad f(x^-) = \inf_{x \in \X} f(x)$$


\Title{Partielle Ableitungen}

Eine partielle Ableitung der Funktion $f: \X \subset \R^n \to \R$, wobei $\X$ offen ist, erhält man, indem man 
alle Variablen bis auf eine als konstant betrachtet und dann nach der einen Variable ableitet.

$$\pdv{f}{x_{0, j}} = \lim_{h \to 0} \frac{f(x_{0, 1}, \dots, x_{0, j} + h, \dots, x_{0, n}) - f(x_0)}{h}$$

Wenn $f: \R^n \to \R^m$ für $x_0 \in \R^n$ so ist
\begin{align*}
  \pdv{f(x_0)}{x_j} := \begin{pmatrix}
    \pdv*{f_1(x_0)}{x_j}\\
    \vdots\\
    \pdv*{f_m(x_0)}{x_j}\\
  \end{pmatrix}
\end{align*}
Eigenschaften von partiellen Ableitungen sind (angenommen die partielle Ableitung von $f, g$ existiert für $x_j$)
\begin{enumerate}[label=(\arabic*)]
  \item $\partial_j(f + g) = \partial_j (f) + \partial_j g$
  \item $\partial_j(f \cdot g) = \partial_j (f) \cdot g + \partial_j (g) \cdot f)$
  \item wenn $g \neq 0$: $\partial_j(f / g) = \frac{\partial_j (f) \cdot g - \partial_j (g) \cdot f}{g^2}$
\end{enumerate}

\textbf{Kettenregel für partielle Ableitungen}.

Sei $f : \R^n \mapsto \R^m$ und $g: \R^m \mapsto \R^k$, wobei $f$ eine Funktion von $n$ Variablen $a_1,...,a_n$ ist 
und $g$ eine Funktion von $m$ Variablen $b_1, ..., b_m$. Nun ist $F(a) = (g \circ f)(a) = g(f(a))$:
$$\partial_{a_i} F = \partial_{b_1} g \cdot \partial_{a_i}f + ... + \partial_{b_m} g \cdot \partial_{a_i}f$$

\Title{Jacobi Matrix}

Sei $f : \X \subset \R^n \to \R^m$ und $\X$ eine offene Menge. Die Jacobi Matrix ist eine Matrix mit $m$ Zeilen und $n$ Spalten:
$$J_f = \Big ( \pdv{f_i}{x_j} \Big ) \underset{\underset{1 \leq j \leq n}{1 \leq i \leq m}}{}$$


\Title{Gradient}

Die Jacobi Matrix einer Funktion $f: \X \subset \R^n \to \R$ ist ein Vektor, der oft mit $\nabla f$
bezeichnet wird. Die geometrische Interpretation ist ein Vektorfeld, definiert durch $\nabla f$, das die 
Richtung und Magnitude des grössten Wachstums von $f$ angibt.


\Title{Divergenz}

Die Divergenz einer Funktion $f$ ist die Spur der Jacobi Matrix von $f$.
$$\text{div}(f)(x_0) = \text{Tr}(J_f(x_0)) = \Sum_{i=0} (J_f)_{i,i}$$


\Title{Differenzierbarkeit}

Sei $\X \subset \R^n \to \R^m$ eine Funktion, $x_0 \in \X$ mit $\X$ offen. Wir sagen $f$ ist differenzierbar 
an der Stelle $x_0$ falls eine lineare Abbildung $L: \R^n \to \R^m$ (d.h. eine $m \times n$ Matrix $L$) 
existiert so dass $\forall x_0 + v \in \X$:

$$f(x_0 + v) = f(x_0) + L(v) + R(x_0, v)$$
$$\lim_{v \to 0} \frac{||R(x_0, v)||}{||v||} = 0$$

Wir schreiben $df(x_0) = L$. Wenn $f$ für alle $x_0 \in \X$ differenzierbar ist, so ist $f$ auf ganz $\X$ differenzierbar. \medskip

Seien $f, g$ differenzierbar im Punkt $x_0 \in \X$ so gilt:
\begin{enumerate}[label=(\arabic*)]
  \item $f$ ist stetig im Punkt $x_0$
  \item $f$ hat alle partiellen Ableitungen am Punkt $x_0$ und die Matrix die die lineare Abbildung 
  $df(x_0): x \mapsto Ax$ repräsentiert ist die Jacobi Matrix von $f$ am Punkt $x_0$, 
  d.h. $A = J_f(x_0)$ (nur $\Rightarrow$)
  \item $d(f + g)(x_0) = df(x_0) + dg(x_0)$
  \item Wenn $m = 1$ so ist $f \cdot g$ differenzierbar und falls $g \neq 0$ $f / g$ auch.
  \item Wenn $f: \X \to Y, g: Y \to \R^m$ beide differenzierbar sind, so gilt 
  $d(g \circ f)(x_0) = dg(f(x_0)) \cdot df(x_0)$. Weiter ist $J_{g \circ f}(x_0) = J_g(f(x_0)) \cdot J_f(x_0)$.
\end{enumerate}

Zuletzt gilt:
$$ \text{$\forall i, \exists\partial_i$ und $\partial_i$ stetig} \Rightarrow f \text{ ist differenzierbar mit } df = J_f$$

Um zu beweisen, dass eine Funktion im Punkt $x_0$ differenzierbar ist, zeigen wir dass:
$$\lim_{x \to x_0} \frac{|f(x) - f(x_0)|}{||x - x_0||} = 0$$

\Title{Tangentialraum}

Der Tangentialraum eines Graphen $f$ am Punkt $x_0$ ist bestimmt durch den Graphen $g(x) = f(x_0) + df(x_0)(x-x_0)$.


\Title{Richtungsableitungen}

Die Richtungsableitung von $f: \X \to \R^m$, wobei $\X$ offen und $x_0 \in \X$ ist, in die Richtung $v \in \R^n$, ist definiert als:
$$w = \lim_{t \to 0} \frac{f(x_0 + t \cdot v) - f(x_0)}{t} = df(x_0) \cdot (v)$$
Falls $v = e_i$ so gilt $w = \partial_i f(x_0)$.


\Title{Höhere Ableitungen}

Wir brauchen Ableitungen von höhere Ordnung für Taylor Polynomen und um Maximas und Minimas
zu bestimmen. \medskip

Sei $\X \subset \R^n$ offen, $f: \X \mapsto \R^m$. Wir sagen $f$ ist eine Funktion der Klasse $C^1$ falls $f$ 
auf $\X$ differenzierbar ist und alle partiellen Ableitungen stetig sind. Wir sagen $f \in C^k$ für $k \geq 2$
wenn $f$ differenzierbar ist und für alle partiellen Ableitungen gilt $\partial_{x_i} f : \in C^{k-1}$. 
Weiter, $f$ ist glatt oder in $C^\infty$ falls $f \in C^k, \; \forall k$. \medskip


%-------------------------
\vfill\null
\columnbreak
%-------------------------


Partielle Ableitungen (bis zur Ordnung $k$) auf einer offenen Menge $\X$ sind kommutativ. Beispiel:
$$\pdv{^2 f}{x_i \partial x_j} = \pdv{^2 f}{ x_j \partial x_i}$$

Polynome sind in $C^\infty (\R^n, \R)$.


\Title{Hesse-Matrix}

Sei $f : \R^n \mapsto \R$. Die Hesse-Matrix ist eine $n \times n$ symmetrische Matrix, 
die die zweite Ableitung definiert.
$$\text{Hess}_f(x_0) := \Big( \pdv{^2 f(x_0)}{x_i \partial x_j} \Big)_{1 \leq i, j \leq n}$$


\Title{Taylor Polynome}

Sei $k \geq 1$ und $f: \X \mapsto \R$ ist eine Funktion der Klasse $C^k$ auf $\X$, sei $x_0 \in \X$ fix. Das $k$-te Taylor Polynom
von $f$ am Punkt $x_0$, ausgewertet am Punkt $x$ ist das Polynom in $n$ Variablen vom Grad $\leq k$:
\begin{align*}
  T_k f(y; x_0) &= f(x_0) + \sum_{i=1}^n \partial_i f(x_0) \cdot y_i + \dots \\
  &+ \sum_{m_1 + \dots + m_n = k} \frac{1}{m_1! \dots m_n!} \pdv{^k f(x_0) }{x_1^{m_1} \dots \partial x_n^{m_n}} \cdot y_1^{m_1} \dots y_n^{m_n}
\end{align*}

Wobei die letzte Summe über alle Tupel von $n$ nicht-negativen ganzen Zahlen geht, wobei die Summe $k$ ist und $y = x - x_0$. Wir können die Notation 
vereinfachen, indem wir für den Multi-Index $m = (m_1, \dots , m_n)$, folgende Notationen einführen:
\begin{align*}
    m!  &= m_1! \cdots m_n! \\
    |m| &= m_1 + ... + m_n \\
    y^m &= y_1^{m_1} \cdots y_n^{m_n} \\
    \partial_x^m f &= \frac{\partial^{|m|} f}{\partial x_1^{m_1} \cdots \partial x_n^{m_n}}
\end{align*}
Nun erhalten wir folgende Notation:
$$T_k f(y; x_0) = \sum_{|m| \leq k} \frac{1}{m!}\partial_x^m f(x_0) \cdot y^m $$

Zum Schluss, wenn $f \in C^k$ für $x_0 \in \X$, so haben wir
$$f(x) = T_k(x - x_0; x_0) + E_k(f, x, x_0)$$
$$\lim_{x \to x_0} \frac{E_k(f, x, x_0)}{\norm*{x - x_0}^k} \to 0$$
\medskip

\textbf{Beispiele:} (k = 1, k = 2)
\begin{align*}
  T_1 f(x; x_0) &:= f(x_0) + \nabla f(x_0) \cdot (x - x_0) \\
  T_2 f(x; x_0) &:= T_1 f(x; x_0) + \frac{1}{2} \cdot (x - x_0)^\top \cdot \text{Hess}_f(x_0) \cdot (x - x_0)
\end{align*}


\Title{Kritische Punkte}

Ein Punkt $x_0 \in \X$ wofür $\nabla f(x_0) = 0$ ist ein \textbf{kritischer Punkt}. Falls $f$ in der Klasse 
$C^2$ ist, so ist ein Punkt nicht-degeneriert, falls $\text{det(Hess}_f(x_0)) \neq 0$


%-------------------------
\vfill\null
\columnbreak
%-------------------------


\Title{Extrema}

Wenn $f: \X \mapsto \R$ differenzierbar ist auf dem Inneren von $\X$, wobei $\X$ kompakt ist, so existieren
globale Extremalstellen von $f$. Diese sind entweder kritische Punkte von $f$ oder befinden sich auf dem Rand
von $\X$.

\Title{Lokale Extrema}  

Sei $f: \X \subset \R^n \mapsto \R$ differenzierbar und $ \X$ eine offene Menge. Wir sagen $x_0 \in \X$ ist ein
\textbf{lokales Maximum (Minimum)} falls wir eine Umgebung 
$B_r(x_0) = \{ x \in \R^n\;|\; \norm*{x - x_0} < r \} \subset \X$
finden können mit:
$$\forall x \in B_r(x_0) \quad f(x) \leq (\geq) f(x_0)$$
Alternativ wird auch von einem Hyperkubus $C(x_0, r)$ gesprochen. Weiter gilt:
$$x_0 \in \X \text{ ist ein lokales Extrema } \Rightarrow \nabla f(x_0) = 0$$

Falls ein kritischer Punkt weder ein lokales Minimum, noch ein lokales Maximum ist, so nennen wir ihn einen
\textbf{Sattelpunkt}.


\Title{Globale Extrema}

Sei $f: K \mapsto \R$ mit $K$ kompakt, dann existiert ein globales Extrema von $f$ und es ist entweder ein kritischer Punkt 
oder am Rand von $K$. Um solch ein Extrema zu bestimmen, teilen wir $K$ in sein Inneres $\X$ und seinen Rand $B$ auf. 
Danach gehen wir wie folgt vor:
\begin{align}
  \text{Bestimmen der kritischen Punkte von } \X \\
  \text{Bestimmen der Minimas und Maximas von } f_{|B}
\end{align}
Beim überprüfen von $X$ kann wie zuvor vorgegangen werden. Das Überprüfen des Randes ist ein Problem der Art $\R \mapsto \R$ 
und kann mit Wissen aus Analysis I gelöst werden.


\Title{Testen von kritischen Punkten}

Sei $f:\X \subseteq \R^n \mapsto \R$, $\X$ offen und $f \in C^2$. Sei $x_0$ ein nicht-degenerierter
kritischer Punkt von $f$. So gilt
\begin{enumerate}[label=(\arabic*)]
  \item Wenn $\text{Hess}_f(x_0)$ pos. def. so ist $x_0$ ein lokales Minimum
  \item Wenn $\text{Hess}_f(x_0)$ neg. def. so ist $x_0$ ein lokales Maximum
  \item Wenn $\text{Hess}_f(x_0)$ ist Indefinit so ist $x_0$ ein Sattelpunkt
\end{enumerate}
Dies können wir nicht machen, falls $x_0$ ein degenerierter kritischer Punkt ist ($\text{det(Hess}_f(x_0)) 
= 0$). In diesem Fall müssen wir einzeln Überprüfen.


\Title{Definit}

Für symmetrische Matrizen haben wir folgende Definition:
\begin{enumerate}[label=(\arabic*)]
  \item Positiv Definit: Alle Eigenwerte sind positiv
  \item Negativ Definit: Alle Eigenwerte sind negativ
  \item Indefinit: Positive und Negative Eigenwerte
\end{enumerate}
Eigenwerte können mit dem charakteristischen Polynom gefunden werden:
\begin{align*}
  \text{det} \Bigg(
  \begin{pmatrix}
    a & b\\
    c & d
  \end{pmatrix}
  -
  \begin{pmatrix}
    \lambda & 0\\
    0 & \lambda
  \end{pmatrix}
  \Bigg)
  &\Rightarrow ad - (a + d) \lambda + \lambda^2 - bc = 0
\end{align*}

Für nicht symmetrische Matrizen, müssen wir testen, ob für jeden Vektor $v$ gilt: $v^\top A v > 0$ (bzw. $<0$). \medskip

\textbf{Marlene Trick}

Sei $M$ eine symmetrische Matrix und $A_k$ die $k \times k$ Blockmatrix der oberen, linken
Ecke. Mit $\triangle_k$ bezeichnen wir die Determinante von $A_k$. Nun ist $M$ pos. def.
wenn $\triangle_k > 0$ für alle $k$, weiter ist $M$ neg. def. wenn $(-1)^k \triangle_k > 0$
für alle $k$.

Für $2$ Dimensionen haben wir
\begin{align*}
  \text{det}
  \begin{pmatrix}
    a & b\\
    c & d
  \end{pmatrix}
  = a \cdot d - c \cdot b 
\end{align*}


%-------------------------
\vfill\null
\columnbreak
%-------------------------


Für $3$ Dimensionen haben wir
\begin{align*}
  \text{det}
  \begin{pmatrix}
    a & b & c\\
    d & e & f\\
    g & h & i
  \end{pmatrix}
  = a \cdot \text{det}
  \begin{pmatrix}
    e & f\\
    h & i
  \end{pmatrix}
  - b \cdot \text{det}
  \begin{pmatrix}
    d & f\\
    g & i
  \end{pmatrix}
  + c \cdot \text{det}
  \begin{pmatrix}
    d & e\\
    g & h
  \end{pmatrix}
\end{align*}


\Title{Max / Min mit Nebenbedingungen}

Wir wollen Minimas und Maximas einer Funktion $f: \X \mapsto \R$, mit einer Nebenbedingung 
$Y = \{x \in \X | g(x) = 0, g:\X \mapsto \R\}$ bestimmen. Dafür können wir die Methode der
Lagrange-Multiplikatoren verwenden.

\begin{tcolorbox}
    Sei $\X \subset \R^n$ offen und $f,g : \X \mapsto \R$ eine Funktion der Klasse $C^1$. Sei 
    $Y = \{ x \in \X | g(x) = 0 \}$, nun ist $x_0 \in Y$ ein lokales Extrema der Nebenbedingung $f_{|Y}$, wenn
    $$\exists \delta > 0, \; f(y) \leq (\geq) f(x_0), \; \; \; \forall y \in B_\delta(x_0) \cap Y$$
\end{tcolorbox}

\begin{tcolorbox}
    Wenn $x_0 \in Y$ eine lokale Extremalstelle von $f_{|Y}$ ist, so ist entweder $\nabla g(x_0) = 0$ oder 
    $\exists \lambda \in \R, \; \nabla f(x_0) = \lambda \cdot \nabla g(x_0)$. Wobei $\lambda$ der
    \textbf{Lagrange-Multiplikator} ist.
\end{tcolorbox}